Three complementary RNA sequencing (RNAseq) platforms were used to profile the transcriptome of \syn~and \syna, each capturing a different facet of transcript structure and abundance.
Illumina short-read sequencing provided the quantitative backbone, PacBio Iso-Seq\textsuperscript{TM} resolved full-length isoform structure, and Oxford Nanopore Technologies (ONT) direct-RNA sequencing provided native, amplification-free long reads.

Illumina-based sequencing is currently the most common technique, owing to its robustness, cost, and throughput.
Illumina instruments generate billions of reads per flow cell, and the chemistry has been refined considerably since the technique was introduced~\cite{mortazavi_mapping_2008}.
Library preparation requires several intermediary steps prior to sequencing: (i) fragmentation of RNA, (ii) conversion of isolated mRNA fragments to double-stranded complementary DNA (cDNA) by reverse transcription, (iii) adapter ligation, (iv) PCR amplification, and (v) size selection~\cite{stark_rna_2019}.
Fragments are sequenced from one end (single-end) or both ends (paired-end), typically with read lengths of tens of nucleotides depending on the instrument.
The reads are short but exceed 99\% base accuracy, particularly with the latest chemistries and instrumentation.

Pacific Biosciences (PacBio) sequencing requires similar library-preparation steps from mRNA to cDNA but sequences full-length cDNA through the Iso-Seq\textsuperscript{TM} protocol, which incorporates SMRTbell adapters for circular consensus sequencing (CCS, also known as HiFi sequencing)~\cite{wenger_accurate_2019}.
Improved instrumentation such as the PacBio Revio and concatemerization techniques such as Kinnex yield accurate full-length isoforms of up to 20 kb, although longer transcripts may be truncated or lost when size selection is applied.
PacBio requires substantial RNA input, so PCR amplification is commonly used during library preparation, and the platform carries a known bias toward smaller fragments that can complicate expression quantification~\cite{rhoads_pacbio_2015}.

Both Illumina and PacBio Iso-Seq convert RNA to cDNA and rely on PCR amplification and size selection, and short reads alone cannot capture the full complexity of the transcripts present in a bacterial cell.
ONT instead captures long transcripts directly by passing native RNA through a membrane-embedded nanopore (\threeprime~to~\fiveprime) and reading the perturbations of an applied electric current~\cite{garalde_highly_2018}.
Direct-RNA sequencing was included for its ability to read full-length native transcripts without conversion to cDNA, which can perturb the native molecule and lose information~\cite{soneson_comprehensive_2019}.
Beyond quantification, native long reads resolve gene co-expression, antisense and intergenic transcription, and, in principle, common base modifications such as N6-methyladenosine~\cite{liu_accurate_2019}.
% The three RNAseq approaches are further reviewed in~\cite{Stark2019,Satam2023}.
